\chapter{Teoría de la especialización del grupo fundamental}
\label{X}
\marginpar{261}
En la presente exposición nos limitamos al estudio del grupo fundamental de
las fibras geométricas en un morfismo \emph{propio}, \ie del grupo fundamental
de un esquema algebraico \emph{propio} variable. En una exposición posterior
generalizaremos la técnica empleada a los recubrimientos étale
\emph{moderadamente ramificados} «en el infinito». Esto nos dará, por ejemplo,
una solución del «problema de los tres puntos» en el caso de los recubrimientos
galoisianos de orden primo con la característica (\ie una determinación de los
recubrimientos galoisianos de la recta proyectiva~$\PP_k^1$, ramificados a lo
sumo en tres puntos dados y moderadamente ramificados en esos puntos), así como
de sus variantes evidentes.

\section[La sucesión exacta de homotopía para un morfismo propio]
{La sucesión exacta de homotopía para un morfismo propio y separable}
\label{X.1}

\begin{definition}
\label{X.1.1}
Un preesquema~$X$ sobre un cuerpo~$k$ se dice \emph{separable}, o
\emph{separable sobre~$k$}, si para toda extensión~$K$ de~$k$,
$X\otimes_k K$ es reducido. Si~$f\colon X\to Y$ es un morfismo de
preesquemas, se dice que~$f$ es \emph{separable}, o que~$X$ \emph{es separable
sobre~$Y$}, si~$X$ es plano sobre~$Y$ y si, para todo~$y\in Y$, la fibra
$X\otimes_Y \kres(y)$ es separable sobre~$\kres(y)$.
\end{definition}

Si~$X$ es un preesquema sobre un cuerpo~$k$, decir que es separable significa
también que es \emph{reducido} y que los cuerpos~$\kres(x)$, para~$x$ punto
genérico de una componente irreducible de~$X$, son extensiones separables
de~$k$. Si~$k$ es perfecto, equivale, pues, a decir que~$X$ es reducido.
Obsérvese que, si~$X$ es separable sobre~$Y$, entonces, para todo cambio de
base~$Y'\to Y$,~$X'=X\times_Y Y'$ es separable sobre~$Y'$. Puede probarse
también, bajo hipótesis de finitud convenientes, que el compuesto de dos
morfismos separables es
\marginpar{262}
un morfismo separable. Solo lo necesitaremos en la forma siguiente:
\emph{si~$X$ es separable sobre~$Y$ y~$X'$ es étale sobre~$X$, entonces~$X'$
es separable sobre~$Y$.} Es, en efecto, una consecuencia inmediata de las
definiciones y de~\Ref{I.9.2}. Por otra parte, la hipótesis «morfismo
separable» nos servirá mediante la proposición siguiente:

\begin{proposition}
\label{X.1.2}
Sea~$f\colon X\to Y$ un morfismo propio y separable, con~$Y$ localmente
noetheriano, y consideremos su factorización de \emph{Stein}
$X\lto{f'}Y'\dpl\to Y$, donde~$f_*'(\cal{O}_X)=\cal{O}_{Y'}$, siendo~$Y'$
finito sobre~$Y$ e isomorfo al espectro del álgebra~$f_*(\cal{O}_X)$. Entonces
~$Y'$ es un \emph{recubrimiento étale} de~$Y$.
\end{proposition}

Esta proposición figurará en (EGA~III 7)\footnote{\Cf EGA~III 7.8.10 (i)}.
Indiquemos el principio de la demostración. Se reduce fácilmente al caso en que
~$Y$ es el espectro de un anillo local completo~$A$ y, haciendo además una
extensión finita plana conveniente de este último (correspondiente a una
extensión residual conveniente), puede suponerse que las componentes conexas
de la fibra del punto cerrado~$y$ son geométricamente conexas, lo que significa
también que~$\H^0(X_y,\cal{O}_{X_y})$ se descompone en un producto de cuerpos
idénticos a~$k=\kres(y)$. Suponiendo entonces~$X$ conexo, como es lícito, se
tendrá~$\H^0(X_y,\cal{O}_{X_y})=k$, por lo que el homomorfismo
$A\to \H^0(X_y,\cal{O}_{X_y})$ es \emph{sobreyectivo}. Se concluye, mediante
una proposición general (de tipo Künneth), que~$f_*(\cal{O}_X)$ está definido
por un módulo~$B$ sobre~$A$ que es libre sobre~$A$, y que
$B/\goth{m} B\to \H^0(X_y,\cal{O}_{X_y})=k$ es biyectivo. Por tanto, en este
caso~$B$ es un álgebra étale sobre~$A$, lo cual termina la demostración.

\begin{theoreme}
\label{X.1.3}
Sea~$f\colon X\to Y$ un morfismo propio y separable, con~$Y$ localmente
noetheriano y conexo, y supongamos que~$f_*(\cal{O}_X)=\cal{O}_Y$ (lo que
implica que las fibras de~$X$ sobre~$Y$ son geométricamente conexas, y
recíprocamente gracias a~\Ref{X.1.2}). Sean~$y$ un punto de~$Y$,
$\overline{\kres(y)}$ una clausura algebraica de~$\kres(y)$ y
$\overline X_y=X_y\otimes_{\kres(y)}\overline{\kres(y)}$. Sean, por último,
~$X'$ un recubrimiento étale \emph{conexo} de~$X$ y
$\overline X_y'=X_y'\otimes_{\kres(y)}\overline{\kres(y)}$. Para que exista
un recubrimiento étale~$Y'$ de~$Y$ y un~$X$-isomorfismo
\marginpar{263}
$X'\isomto X\times_Y Y'$, es necesario y suficiente que~$\overline{X'}_y$
admita una sección sobre~$\overline X_y$.
\end{theoreme}

Poniendo~$Y'=\Spec(h_*(\cal{O}_{X'}))$ (donde~$h\colon X'\to Y$ es el
morfismo compuesto~$X'\to X\to Y$), basta probar que el~$Y$-morfismo canónico
$$
X'\to X\times_Y Y'
$$
es un \emph{isomorfismo} y que~$Y'$ es étale sobre~$Y$. Ahora bien, ya sabemos
por~\Ref{X.1.2} que~$Y'$ es étale sobre~$Y$; por tanto,
~$X\times_Y Y'$ es étale sobre~$X$, y por ello el morfismo
~$X'\to X\times_Y Y'$ es también étale (\Ref{I.4.8}). Además,~$Y'$ es
conexo como imagen de~$X'$, que lo es; por tanto,~$X\times_Y Y'$ es conexo,
pues~$X$ tiene fibras conexas sobre~$Y$ (\Ref{IX.3.4} y
\Ref{V.6.9}~(iii)). Así, para probar que~$X'\to X\times_Y Y'$ es un
isomorfismo, basta ver que su grado de proyección en \emph{un} punto de
~$X\times_Y Y'$ es igual a~1. Esto resulta fácilmente de la hipótesis de que
~$\overline{X'}_y$ admite una sección sobre~$\overline X_y$, bien utilizando
\Ref{IX.6.6}, bien, más simplemente, observando que basta probar la
existencia de tal punto en~$X\times_Y Y'$ después del cambio de base
\ifthenelse{\boolean{orig}}
{$\Spec(\overline k)\to Y$}
{$\Spec(\overline{\kres(y)})\to Y$}, donde esto es evidente.
Esto termina la demostración de~\Ref{X.1.3}.

Teniendo en cuenta~\Ref{IX.3.4} y el diccionario~\Ref{V.6.9} y
\Ref{V.6.11}, puede ponerse~\Ref{X.1.3} en la forma equivalente siguiente:

\begin{corollaire}
\label{X.1.4}
Con las notaciones precedentes para~$f\colon X\to Y$ y~$\overline X_y$, sea
~$\bar{a}$ un punto geométrico de~$\overline X_y$,~$a$ su imagen en~$X$ y
~$b$ su imagen en~$Y$. Entonces la sucesión siguiente de homomorfismos de
grupos es \emph{exacta}:
$$
\pi_1(\overline X_y,\bar{a})\to \pi_1(X,a)\to \pi_1(Y,b)\to e.
$$
\end{corollaire}

\begin{remarques}
\label{X.1.5}
Obsérvese que la demostración de~\Ref{X.1.3} hace intervenir de manera esencial
~\Ref{X.1.2} y, por ello, el «primer teorema de comparación» en geometría
algebraico-formal. En cambio, la teoría del descenso de la exposición~\Ref{IX}
\marginpar{264}
solo ha intervenido mediante~\Ref{IX.3.4}, cuya demostración directa es
fácil en el caso de un morfismo \emph{propio}~$f\colon X\to Y$ tal que
~$f_*(\cal{O}_X)=\cal{O}_Y$. Sea, en efecto,~$Y'$ étale sobre~$Y$ y
supongamos que~$X'=X\times_Y Y'$ sea unión disjunta de dos abiertos no vacíos;
probemos que ocurre lo mismo con~$Y'$. En efecto, se tendrá
~$Y'=\Spec(\cal{A})$, y por tanto~$X'=\Spec(\cal{B})$, con
~$\cal{B}=\cal{A}\otimes_{\cal{O}_Y}\cal{O}_X$; la descomposición de~$X'$ en
unión disjunta corresponde a una descomposición de~$\cal{B}$ en producto de
dos álgebras no nulas~$\cal{B}_1$ y~$\cal{B}_2$. Como
~$f_*(\cal{O}_X)=\cal{O}_Y$, se concluye fácilmente que
~$f_*(\cal{B})=\cal{A}$; por tanto,~$\cal{A}$ será producto de dos álgebras
(también no nulas, pues sus secciones unidad no son nulas)
~$f_*(\cal{B}_1)$ y~$f_*(\cal{B}_2)$, c.q.d.
\end{remarques}
